\documentclass[12pt]{report}
\usepackage[utf8]{inputenc}
\usepackage[top = 1.5cm, bottom = 1.5cm, left = 1.0cm, right = 1.0cm]{geometry}
\usepackage[backend=biber,style=numeric,sorting=none]{biblatex}
\usepackage[dvipsnames]{xcolor}
\usepackage[T1]{fontenc}
\usepackage[french]{babel}
\usepackage[title]{appendix}
\usepackage[most]{tcolorbox}
\usepackage{amsmath}
\usepackage{amssymb}
\usepackage{amsfonts}
\usepackage{adjustbox}
\usepackage{caption}
\usepackage{csquotes}
\usepackage{fancyhdr}
\usepackage{graphicx}
\usepackage{hyperref}
\usepackage{lmodern}
\usepackage{listings}
\usepackage{lastpage}
\usepackage{longtable}
\usepackage{minted}
\usepackage{mathtools}
\usepackage{ragged2e}
\usepackage{subcaption}
\usepackage{tocloft,url}
\usepackage{threeparttable}
\usepackage{url}
\usepackage{verbatim}
\usepackage{xcolor}

%----------------------------------------------------------------------------------
%									  COMMANDES
%----------------------------------------------------------------------------------
\DeclarePairedDelimiter\abs{\lvert}{\rvert}%
\DeclarePairedDelimiter\norm{\lVert}{\rVert}%

\usemintedstyle{fruity}
\definecolor{bg}{HTML}{282828}

\setlength{\headheight}{15.35403pt}
\setlength{\parskip}{0.5em}

\setcounter{secnumdepth}{3}

\newlistof{links}{lks}{Liste des liens}
\newcommand\externallink[2]{%
\refstepcounter{links}%
\footnote{#1\url{#2}}%
\addcontentsline{lks}{links}{%
\protect\numberline{\thelinks}%
\protect{\url{#2}}}%
}

\hypersetup{
	colorlinks=true,
	linkcolor=blue,
	filecolor=magenta,      
	urlcolor=blue,
	citecolor=darkgray,
	pdftitle={Rapport de stage}
}

\newcommand\NB[1][0.3]{N\kern-#1em{B}}

%%Création commande pour insérer image avec nom d figure directement
%\newcommand{nomDeTaCommande}[nombreArguments]{CodeLaTeX}
%\insertImage[position]{image_path}{scale}{Titre_figure}{label}
\newcommand{\insertImage}[5][center]{
	\begin{#1}
		\includegraphics[scale=#3]{#2}
		\captionof{figure}{#4}
		\label{#5}
	\end{#1}
	}
	
	%%Création d'une nouvelle commande pour faire référence à une Figure
	%Exemple : \appelFigure{schema} donne : Figure 1 (en italique)
	\newcommand{\appelFigure}[1]{
		\textit{Figure \ref{#1}}
		}
		
		%%Création d'une nouvelle commande pour créer une barre horizontale
		\newcommand{\HRule}{\rule{\linewidth}{0.5mm}}
		
		\renewcommand{\contentsname}{Sommaire}
		
		\DeclareMathOperator*{\argmin}{arg\,min}
		\DeclareMathOperator*{\argmax}{arg\,max}
		
		%----------------------------------------------------------------------------------
		%									INITIALISATIONS
		%----------------------------------------------------------------------------------
		\hbadness 11000
		\vbadness 11000
		\interlinepenalty 10000
		
		\title{{\huge \bfseries Rapport de TP}\\[0.2cm]
		% \includegraphics[scale=0.75]{../images/sheet.png} \\[0.4cm]
		{\Large Intelligence Artificielle}}
		\author{Théo \textsc{FIGINI}\\L3 Informatique\\Année universitaire 2021-22\\ [1.5cm]
				Syoan \textsc{ODOUHA}\\L3 Informatique\\Année universitaire 2021-22}
		\date{Dimanche 3 Mars 2022}
		\makeatletter
		\setcounter{chapter}{0}
		
		\pagestyle{fancyplain}
		\fancyhf{}
		\lhead{Rapport de TP}
		\rhead{\leftmark}
		\cfoot{Page \thepage /\pageref{LastPage}}
		\renewcommand{\headrulewidth}{0.5pt}
		\renewcommand{\footrulewidth}{0.5pt}
		\renewcommand{\plainheadrulewidth}{0.5pt}
		\renewcommand{\plainfootrulewidth}{0.5pt}
		\renewcommand\labelitemi{--}
		\renewcommand\topfraction{.9}
		\renewcommand\textfraction{0.35}
		\renewcommand\floatpagefraction{0.8}
		
		% \bibliography{biblio}
		
		\counterwithout{figure}{chapter}
		
		
\begin{document}
\begin{titlepage}
	\begin{center}
		\includegraphics[scale=1.80]{./images/iconUA.png}
		\hspace{2cm}
		\includegraphics[scale=0.25]{./images/iconSpike.png} \\[2cm]
		\hspace{2.5cm}

		\HRule \\[0.4cm]
		\@title \\[0.4cm]
		\HRule \\[1cm]

		\@author \\ [1.5cm]

		{\large Organisme d'accueil: \textsl{Université des Antilles}} \\[1.5cm]

		\begin{minipage}{0.7\textwidth}
			\begin{center}
				Enseignant: \\
				\hspace{0.2cm} Vincent \textsc{PAG\'E}
			\end{center}
		\end{minipage}\\[3cm]

		\@date
	\end{center}
\end{titlepage}


\renewcommand{\thefigure}{\arabic{figure}}
\setcounter{figure}{0}

\chapter{Réalisation}

\section{}
% \chapter{Conclusion}

Au cours de ces travaux pratiques, nous avons étudié les différents algorithmes de classification et de régression. Nous avons pu
en déduire la régression n'était pas une méthode adaptée pour traiter ce problème. Les méthodes ayant eu les meilleures performances
en apprentissage sont le KNN et les SVM, en revanche, le réseau de neurones pour la classification a obtenu les meilleures performances
en validation.

Afin d'améliorer ces résultats, nous aurions pu essayer de choisir des caractéristiques plus pertinentes ou améliorer le traitement de nos données.

% \printbibliography
\end{document}
